\begin{table}[htbp]
\raggedright
\caption{Mining Exposure Covariates, Coal Mining Exposed Utilities, Two-Step Upstream Watershed Linkage, 1985--2005}
\label{tab:mining_exposure_covariates_k2}
\small
\begin{tabular}{>{\raggedright\arraybackslash}p{3cm} *{4}{>{\centering\arraybackslash}p{1.15cm}}}
\toprule
\textbf{Variable} & \textbf{Mean} & \textbf{SD} & \textbf{P90} & \textbf{P99} \\
\hline
Number of mines upstream & 1.12 & 2.27 & 4.00 & 11.00 \\
Cumul. upstream coal prod. (10M ST) since 1985 & 0.32 & 1.00 & 0.68 & 5.23 \\
Percentage of coal weight as sulfur & 1.11 & 1.13 & 2.96 & 4.01 \\
\bottomrule
\end{tabular}
\begin{minipage}{\linewidth}
\vspace{4pt}
\footnotesize
\raggedright
\textit{Notes:} Sample of drinking water utilities with a coal mine within two flow steps upstream of their intake and no coal mine colocated with their intake, 1985--2005. Mean, SD, P90, and P99 are reported for the number of coal mines upstream, cumulative coal production upstream since 1985 (10 million short tons), and the percentage of coal weight that is sulfur, for the same sample. Number of observations = Number of utilities $\times$ Number of years. N\,=\,10,782 = 582 utilities $\times$ up to 21 years (1985--2005).
\end{minipage}
\end{table}
